\section{Related Work and Architecture Positioning}
\label{sec:positioning}

This paper extends \textsc{Probingnoise} by layering a new module on top of its existing core rather than replacing that core. The implemented base system is a synchronous, Colab-trainable stack composed of a local temporal encoder, a graph-coordination block with residual skip and layer normalization, a local--global aggregation block with residual skip, and a task head, with residual gating and adaptive scheduling retained as optional extensions. All empirical results reported in this work are produced by this implementation under synchronous execution.

We position the base architecture in relation to a broader asynchronous-computation lineage only at the level of theoretical context and future systems direction, not as a description of the present implementation. Foundational work on parallel and distributed iterative methods established the mathematical basis for computation without strict global synchronization~\citep{bertsekas1989parallel}, and lock-free stochastic optimization later demonstrated that practical large-scale learning can proceed without global coordination~\citep{niu2011hogwild}. More recent work has explored asynchronous message passing as an alternative to standard synchronous graph neural network execution~\citep{faber2022async,faber2024gwac}. These references establish a principled direction for future systems-level development of \textsc{Probingnoise}; in the present paper, however, \textsc{Probingnoise} itself remains strictly synchronous, and no property demonstrated in those works is claimed as a property of the current implementation.

On top of this base, we specify a visual-evidence extension for multimodal prediction. The extension consists of a visual front-end, a visual representation encoder, a bridge module that projects visual representations into the \textsc{Probingnoise} stack, and a prediction head operating over the fused multimodal representation produced by that bridge. It is defined as additive: it consumes the outputs of the existing core without modifying its internal blocks, residual pathways, or training regime. The purpose of the extension is architectural---to define how visual evidence may be incorporated into the residual structure of \textsc{Probingnoise} while preserving the underlying training regime and continuity of the original stack. No claim is made that the current study implements, trains, or empirically validates this extension.

The design of the extension is grounded in existing work on transformer-based image representation~\citep{dosovitskiy2021vit}, vision--language pretraining via contrastive objectives~\citep{radford2021clip}, frozen-encoder bridging between visual and language models~\citep{li2023blip2}, and multimodal chain-of-thought prompting in language models~\citep{zhang2023mcot}. These references establish the provenance of the component choices; no empirical property demonstrated in a cited work is inherited as a property of the present system. Novelty is therefore claimed not at the level of the constituent visual or multimodal components, which are adopted from prior work, but at the level of integration---specifically, how visual representations are introduced into the \textsc{Probingnoise} residual structure without replacing the existing synchronous core. The claim of the present paper is accordingly narrow and architectural: it specifies a conservative extension path from the current \textsc{Probingnoise} system toward multimodal prediction, and defers empirical validation of that extension to future work.
